% ID: SEL-D0010-BC-E05
% ESTADO: cotejada. Segunda pasada visual completa y borrador compilado revisados.
% Fuente: Selectividad/matematicas_aplicadas/EXTRA_TITULAR_MAT-APLIC-CCSS-JULIO-IMPRESO.pdf
% SHA-256: 548f60ee334516c7fa918dba3c28fb3ea5fe753cf76eb05402d7b8acb5a3b560
% Página: 2; bloque C; ejercicio original 5.

En una población, se sabe que el 15\% de las personas padece una determinada
enfermedad. Si la persona está enferma, un test da positivo en el 92\% de los
casos, mientras que si la persona está sana, el test da positivo en el 4\%
de los casos (falso positivo). Se elige una persona al azar de esa población.
\begin{enumerate}[label=\alph*)]
  \item \textbf{(1.25 puntos)} Calcule la probabilidad de que, habiendo dado
  positivo el test, la persona esté enferma.
  \item \textbf{(0.5 puntos)} Calcule la probabilidad de que la persona esté
  enferma y el test salga negativo.
  \item \textbf{(0.75 puntos)} Calcule la probabilidad de que saliendo el
  test negativo, la persona esté enferma.
\end{enumerate}
